\subsection{Latest upstream results: implications and corrections}
\label{tht:upstream}

On 14 November 2026 the shared checkout was fast-forwarded from
\texttt{ff4afe4} to \texttt{d5e4297} of
\href{https://github.com/qedsphere/AAconjecture}{qedsphere/AAconjecture}.
The new material is
\href{https://github.com/qedsphere/AAconjecture/blob/d5e429763b9f0234f5dd02c3fe03925c78b4306c/tasks-1-2-3-summary.md}
{\emph{Tasks 1--3: results, proofs, counterexamples, and open gaps}}.
It contains a Gibbs/DFKO influence programme, a face-constrained
transportation problem, and a Fourier-completely-bounded (fcb) programme.
Its project-level results and conditional reductions are not imported
as newly verified theorems. Several referenced proof files are absent
from this checkout. The upstream file is preserved unchanged; the
following is our independent audit.

\textbf{Subsequent upstream update.} Before publication we also pulled
\texttt{9312412}, adding the
\href{https://github.com/qedsphere/AAconjecture/blob/9312412/tasks-1-2-3-proof-companion.md}
{Tasks 1--3 proof companion}. It confirms uniform Fourier-level
additivity, corrects the block prefactor to $\delta_{\min}^{6}$, and
distinguishes covariance operator norm from absolute row sums. The
corrections below address the original summary; those three repairs
are acknowledged, and the block theorem remains explicitly not
independently revalidated in the companion. The remaining status and
parameter-dependence questions are recorded in the companion reviews.

\subsubsection{What strengthens the structural theory}

\textbf{Quantifiers matter.} The handoff's Gibbs target permits a
prefactor $c(d,\ldots)$ and variance exponent $C(d,\ldots)$ depending
arbitrarily on degree. Even dimension-independent bounds of that form
need not yield a polynomial in $T,1/\epsilon$. Its hypotheses also
include bounded degree, range, and fields. They therefore do not settle
our conjecture uniform over fields and graph size at a fixed Dobrushin
margin. Fixed-degree influence results remain useful probes when their
degree and structural dependence are stated explicitly.

\textbf{Gibbs averaging is not coordinate restriction.} On a zero-field
three-leaf Ising star, with center 0 and unit couplings,
\[
 \mathbb E[X_0\mid X_1,X_2,X_3]
 =c_1(X_1+X_2+X_3)+c_3X_1X_2X_3,
 \qquad
 c_3=\frac{\tanh(3\beta)-3\tanh\beta}{4}
     =-\frac{2\tanh^3\beta}{1+3\tanh^2\beta}.
\]
This is nonzero for every $\beta>0$, including the strict Dobrushin
regime $3\tanh\beta<1$. Conditional projections and heat-bath
evolution can therefore raise degree from one to three. This rules out
the handoff's degree-preservation shortcut. It does not affect our
forest theorem: that proof bounds the queries of the composed quantum
oracle and samples edge sources conditionally, without claiming degree
preservation under Gibbs averaging. In
Proposition~\ref{tht:reveal-prop}, pinning substitutes known coordinates
in $Q$; conditional means are only leaf outputs and need not be
low-degree polynomials of the remaining coordinates.

\textbf{Physical locality does not localize arbitrary query observables.}
Under independent uniform bits in $\{0,1\}$, take $p=x_ix_j$ with $j$
outside a prescribed neighborhood of $i$. Then
\[
 (p(x)-p(x^{\oplus i}))^2=x_j,
 \qquad
 \inf_{\ell\text{ measurable near }i}\mathbb E(x_j-\ell)^2=\tfrac14.
\]
This bounded degree-two example has no interactions at all. Correlation
decay cannot justify localizing every acceptance gradient around its
differentiation site. Geometry should control the input representation
and its conditional access; the quantum observable may couple distant
sites. The handoff's joint-endpoint TV estimate is also weaker than
the full-input conditional mismatch bound and consistent adaptive
posterior oracle required in
Proposition~\ref{prop:threshold-local-posterior}.

\subsubsection{The transportation problem: existence versus a new proof}

The handoff labels the exact face-supported summed-$(b')$ LP open, but
its feasibility already follows from the accepted determinant argument
in the supplied
\href{https://arxiv.org/abs/2611.11626}{Blanc--Docter--Strassle--Tan
exposition, Lemma 3.1 and Section 2.5.4}. That argument produces a
\emph{permutation} $\pi$ of $A$ satisfying the face constraint and,
simultaneously for every coordinate,
\[
 \#\{x\in A:\pi(x)_i=x_i\}
 \le\#\{x\in A:x^{\oplus i}\notin A\}.
\]
Set $M_{xy}=\mathbf1[y=\pi(x)]$. It is doubly stochastic and has the
required face support. Multiplying these inequalities by $w_i>0$ and
summing proves exactly the handoff's weighted LP bound. The remaining
research problem is a \emph{determinant-free proof or construction},
not existence or a counterexample to that LP under the imported theorem.
Birkhoff decomposition alone does not give such a new construction;
Hall feasibility alone does not imply the weighted cost estimate.
No weighted Gibbs-law analogue is thereby proved.

Indeed a literal Gibbs-preserving version already fails for one biased
independent bit. Let $A=\{0,1\}$ and $T(0)=T(1)=\{1\}$. For a law
with $\mu(0)=p\ne1/2$, face support forces any coupling with first
marginal $\mu$ to have masses $\gamma(0,1)=p$ and
$\gamma(1,0)=1-p$. Its second marginal is reversed, so it cannot also
equal $\mu$. Both faces are valid. Even empty faces would force this
swap if one retained the zero-boundary, zero-agreement-cost bound.
This is an obstruction to that literal transport extension, while
independent-field inputs remain full-AA positive. A nonuniform proof
must change the coupling, the cost, or the source representation; the
failure of this transport cannot define a temperature threshold.

\subsubsection{Corrections needed before importing the other claims}

\begin{enumerate}
\item \textbf{Uniform Fourier levels have additive influences.}
If $p=\sum_{k=0}^d p_k$ is its homogeneous Walsh decomposition, then
with our normalized flip convention,
\[
 I_i^U(p)=\sum_{S\ni i}\widehat p(S)^2
         =\sum_{k=1}^d I_i^U(p_k).
\]
The contrary sentence in the handoff is false under uniform input.
What a levelwise fcb argument must still track is the norm of the
selected level, including its ambient degree parameter. Under
correlated laws, Walsh orthogonality need not survive.

\item \textbf{The stated polynomial fcb target is already a uniform-AA
corollary.} The primary definition satisfies
$\|p\|_\infty\le\|p\|_{\mathrm{fcb},d}$; see
\href{https://arxiv.org/abs/2304.06713v2}{Escudero Guti\'errez,
Proposition 2.3}. Hence $\|p\|_{\mathrm{fcb},d}\le1$ is a stronger
restriction than supremum boundedness and puts $p$ within the
uniform-AA theorem imported here. The sharper degree-independent
$\operatorname{Var}^2$ estimate, or a proof exploiting fcb directly,
is a separate target. This supplies no correlated-Gibbs transfer.

\item \textbf{Repair the mean-square comparison witness.}
The handoff's $f=1+n^{-1}\sum_iY_i$, with independent mean-$1/2$
signs $Y_i$, gives a squared-norm ratio tending to $9/4$; that example
alone does not refute an arbitrary finite comparison constant. The
intended claim is proved by the bounded affine polynomial
$M=n^{-1}\sum_iY_i$:
\[
 \mathbb E_U M^2=1/n,\qquad
 \mathbb E_\mu M^2=\tfrac14+\tfrac{3}{4n},\qquad
 \frac{\mathbb E_\mu M^2}{\mathbb E_U M^2}=\frac{n+3}{4}.
\]
Independent fields themselves have full AA. This is a failure of a
comparison method, not of simulation.

\item \textbf{Do not use the reported inverse-atom prefactor.}
The handoff's project-level two-bit influence formula has a factor
$\delta_{\min}^{-6}$. It cannot hold with a positive prefactor uniform
in $\delta_{\min}$: take a nearly-copy two-bit law with fair marginals
and $p(x)=x_1$ on $\{0,1\}^2$. Its variance is $1/4$ and unnormalized
maximum flip influence is one, while the proposed right side diverges
as the least atom probability tends to zero. The source formula needs
its exact proof and parameter dependence; we do not silently repair
its exponent or rely on it. Our full-AA block theorem has its own proof.

\item \textbf{Boolean degree does not imply $d^2\ge n$.}
For the addressing function with $k$ address bits and $2^k$ data bits,
all $n=k+2^k$ variables are relevant, while the exact degree is $d=k+1$.
Expand it as $\sum_a\mathbf1[\text{address}=a]x_a$ to see the upper
bound; the top coefficient involving a fixed data bit is nonzero.
The handoff's blanket Boolean shortcut in Task 2 is therefore unavailable.

\item \textbf{Critical covariance needs a separate argument.}
On a free $b$-ary Ising tree of depth $h$,
$\operatorname{Cov}(X_o,X_v)=\tanh(\beta)^{\operatorname{dist}(o,v)}$.
At $\tanh\beta=1/b$, the root covariance row sum is $h+1$, not a
dimension-uniform constant. Thus that proof route cannot establish the
reported critical endpoint without another argument. This does not
refute a separate influence theorem, and our forest simulation does
not use covariance-row summability. On the square lattice the usual
field-uniform coefficient bound is $4\tanh\beta<1$, reaching
$\operatorname{arctanh}(1/4)$; $\beta<1/4$ is the more conservative
linear row-norm criterion. Neither should be conflated with a
convention-free exact physical or full-AA threshold.
\end{enumerate}

\subsubsection{Keeping the quantum target in view}

Restricting attention to genuine quantum acceptance polynomials still
does not make small correlated flip influences a full-AA obstruction.
For example
\[
 Q(x)=\frac12\left(1+\frac1n\sum_i x_i\right)
\]
is the acceptance probability of a one-query algorithm: choose a
uniform coordinate and accept according to its bit. In ordered
Curie--Weiss, $\operatorname{Var}_\mu Q$ stays bounded away from zero
while $I_i^\mu(Q)=1/(4n^2)$. Nevertheless coordinate sampling
approximates $Q$ on every fixed input with polynomially many queries
in inverse accuracy. Even quantum query realizability leaves the gap
between influence failure and simulation failure intact.

Tree reconstruction is another oracle-model distinction. Its hidden
root is inferred from distant boundary spins. Our forest theorem exposes
\emph{all} spins as query coordinates and remains valid on both sides
of uniqueness and reconstruction at every fixed finite temperature.
Boundary-only lower bounds do not apply to that access model.

The new material sharpens the proof programme without changing the
established full-AA temperature classifications. A general local
high-temperature theorem still needs a degree-uniform,
transcript-stable bridge to acceptance simulation. The candidate
bridges remain polynomial reveal progress and local posterior access.
The handoff's withdrawn MLSI/junta citation, unproved cumulant package,
and unfinished determinant-free construction do not supply those bridges.
